\section{Introduction}
\label{sec:intro}

Semantic segmentation analyzes every pixel and classifies each one into related class. 
A vehicle needs to know where pedestrians are and where the road is. Without pixel-level analysis, this task is difficult. 
The Cityscapes dataset~\cite{cityscapes} is a standard benchmark for this type of task. It includes 5000 photos of German streets. 
Each photo is manually labeled with 19 classes such as road, sidewalk, building.
ERFNet~\cite{erfnet} is a traditional pixel-based approach. 
It uses a CNN (convolutional neural network) with an encoder-decoder architecture, and classifies each pixel independently.
MaskFormer~\cite{maskformer} changed this approach by predicting N binary masks and N class labels, then combining them. 
In this way, the model handles both semantic and panoptic segmentation~\cite{panoptic}. 
Mask2Former~\cite{mask2former} improved MaskFormer by adding masked cross-attention, so the model only looks at the image region where each mask is predicted. 
EoMT~\cite{eomt} simplified this further. It uses only the DINOv2~\cite{dinov2} encoder and removes the pixel decoder entirely.
The problem is that all these models are closed-set. 
They were trained on exactly 19 classes, and during training they did not see anything except what were trained on. 
However, the model must assign objects it has never seen to one of its known classes. 
if the model assigns an animal or fence as road, the vehicle tries to drive through it. This creates safety problems. 
This type of data is called out-of-distribution (OoD) data. 
We use four post-hoc scoring methods (MSP, MaxLogit~\cite{maxlogit}, MaxEntropy~\cite{maxentropy}, RbA~\cite{rba}) with temperature scaling, and evaluate them on benchmarks from SMIYC~\cite{smiyc} and Fishyscapes~\cite{fishyscapes}.
We evaluated EoMT on semantic segmentation and compared two differently trained models (COCO-EoMT vs Cityscapes-EoMT).
Then we fine-tuned the COCO-EoMT on Cityscapes, and the model achieved 78.85\% mIoU. 
